\chapter{Análise e Desenho da Solução}
\label{cap:analise_desenho}

% TODO: Introdução curta — este capítulo parte do domínio do problema (quem são os
% utilizadores e o que fazem) e chega às decisões de arquitetura que guiaram a
% implementação, com foco na camada de backend e na infraestrutura.

%----------------------------------------------------------------------------------------
\section{Domínio do Problema}
\label{sec:domínio_problema}

% Apresentar um glossário dos conceitos centrais da plataforma e um modelo de domínio.
%
% CONCEITOS A DEFINIR:
%   - Instituição: entidade que aloja idosos e subscreve o serviço CompyAI.
%   - Senior: utente da instituição; interage com o assistente exclusivamente por voz.
%   - Caregiver: profissional de saúde responsável por um conjunto de seniores;
%     consulta o dashboard de monitorização.
%   - Administrador: gere a instituição na plataforma (contas, perfis, configurações).
%   - Conversa: sessão de interação por voz entre o senior e o assistente de IA.
%   - Sumário: artefacto estruturado gerado automaticamente após cada conversa,
%     com tópicos abordados, tom emocional e pontos-chave.
%   - Medicação: registo dos medicamentos do senior com horários e confirmações de toma.
%
% MODELO DE DOMÍNIO:
%   Diagrama UML de classes simplificado com as relações entre Instituição,
%   Administrador, Caregiver, Senior, Conversa e Sumário.
%   Destacar que um Caregiver gere múltiplos Seniores e que cada Senior tem um
%   histórico de Conversas com os respetivos Sumários — esta estrutura determina
%   diretamente o modelo de dados relacional.

%----------------------------------------------------------------------------------------
\section{Requisitos Funcionais e Não Funcionais}
\label{sec:requisitos}

% REQUISITOS FUNCIONAIS (por perfil de utilizador):
%
% Senior:
%   RF01 — Iniciar e terminar uma conversa por voz com o assistente.
%   RF02 — Receber respostas adaptadas ao seu perfil (tom, vocabulário, duração).
%   RF03 — Ser lembrado de tomar medicação nos horários definidos.
%
% Caregiver:
%   RF04 — Consultar o historial de conversas dos seniores sob a sua responsabilidade.
%   RF05 — Visualizar sumários estruturados de cada conversa.
%   RF06 — Receber alertas quando o sistema deteta sinais de mal-estar.
%
% Administrador:
%   RF07 — Criar e gerir contas de caregiver e senior associadas à sua instituição.
%   RF08 — Configurar o perfil do senior (notas clínicas, medicação, preferências de voz)
%          que serve de contexto ao modelo de IA.
%   RF09 — Consultar estatísticas de utilização da plataforma.
%
% REQUISITOS NÃO FUNCIONAIS:
%   RNF01 — Segurança: autenticação por JWT, autorização por papel (RBAC), TLS em todas
%            as comunicações externas, isolamento do microserviço de IA da internet pública.
%   RNF02 — Disponibilidade: serviço operacional 24/7, dado o contexto de cuidados contínuos.
%   RNF03 — Latência: a API de backend deve responder em menos de 200ms nos endpoints
%            críticos; o pipeline de voz ponta-a-ponta deve produzir a primeira resposta
%            audível em menos de 3 segundos.
%   RNF04 — Privacidade: conformidade com o RGPD no tratamento de dados de saúde;
%            os dados dos utentes não devem sair da infraestrutura controlada pela empresa.
%   RNF05 — Manutenibilidade: arquitetura em camadas com cobertura de testes automatizados,
%            migrações de base de dados versionadas e processo de deploy reprodutível.

%----------------------------------------------------------------------------------------
\section{Arquitetura da Solução}
\label{sec:desenho}

% VISÃO GERAL DO SISTEMA:
%   Diagrama de componentes com os três serviços (interface web, backend, microserviço de IA),
%   a base de dados PostgreSQL e os fornecedores externos de IA.
%   Enfatizar que o microserviço de IA não está exposto à internet — apenas o backend
%   comunica com ele numa rede interna Docker, o que limita a superfície de ataque.
%
% ARQUITETURA INTERNA DO BACKEND (foco principal):
%   Diagrama de módulos mostrando a separação explícita em três camadas:
%     - Controladores (HTTP): recebem e validam os pedidos, aplicam guards de autenticação
%       e autorização, e delegam na camada de serviços.
%     - Serviços (lógica de negócio): contêm as regras do domínio, independentes de
%       frameworks e de infraestrutura — testáveis de forma isolada com mocks.
%     - Repositórios (infraestrutura): encapsulam o acesso à base de dados via Drizzle ORM;
%       a camada de serviços nunca conhece os detalhes de persistência.
%   Justificar esta separação: cada camada pode evoluir e ser testada independentemente;
%   substituir o ORM ou a base de dados não exige alterações na lógica de negócio.
%
% MODELO DE DADOS:
%   Diagrama entidade-relação com as tabelas principais:
%     users, institutions, admins, caregivers, seniors, conversations,
%     conversation_summaries, medications.
%   Decisões de modelação a justificar:
%     - A tabela users centraliza autenticação; admins/caregivers/seniors são extensões
%       com campos específicos de cada perfil (padrão de herança em tabelas separadas).
%     - A tabela conversation_summary é separada de conversations porque o sumário é
%       gerado de forma assíncrona no final da sessão — a conversa existe e é persistida
%       em tempo real, o sumário é produzido depois.
%     - Todas as entidades de negócio têm uma foreign key para institutions, garantindo
%       isolamento de dados entre instituições diferentes (multi-tenancy por linha).
%
% ARQUITETURA DE INFRAESTRUTURA:
%   Diagrama de deployment mostrando os quatro contentores Docker (frontend, backend,
%   serviço de IA, PostgreSQL), a rede interna que os interliga, a instância EC2,
%   o proxy Nginx e a terminação TLS.
%   Mencionar a separação entre rede pública (apenas o Nginx está exposto) e rede
%   interna Docker (backend, serviço de IA e base de dados).
